\section{Local structure}\label{locstr}
In this section, we study the local structure of CAVBs and CLAs with constant real rank. Our main tools for providing these local structures come from the techniques developed in \cite{bursztyn2019splitting}. In~Appendix~\ref{preliminarieslocstr}, we recall the properties of normal vector bundles and Euler-like vector fields, which play a key role in this section.
\subsection{Local description of involutive CAVBs}%\label{locstrcavb}
Involutive RAVBs inherit the structure of an almost Lie algebroid \cite[Proposition 3.17]{bursztyn2019splitting}. For~involutive CAVBs, we have the same result and the proof is an adaptation of the real case to the complex case.
\begin{Proposition}
A real or complex anchored vector bundle $(A,\rho)$ is involutive if and only if there exists a bracket on $\Gamma(A)$ making $A$ into a real or complex almost Lie algebroid, respectively.
\end{Proposition}

Let $(A,\rho)$ be an involutive CAVB with constant real rank. Consider a bracket $[\cdot,\cdot]$ that makes~$(A,\rho)$ into an almost Lie algebroid. Then, we have an operator $D_{\sigma}\colon \Gamma(A)\to \Gamma(A)$, given by $D_{\sigma}\tau=[\sigma,\tau]$. As mentioned in \cite{bischoff2020deformation}, we have the map
\begin{gather*}
\widetilde{\rho}\colon \ \Gamma (A^{\real}) \to \mathfrak{aut}(A,\rho),\qquad
\widetilde{\rho}(\sigma)=(D_{\sigma}, \rho(\sigma))
\end{gather*}
that fits into the following diagram:
\begin{equation}{\label{lifting_aut}}
\xymatrix{
& \mathfrak{aut}(A,\rho)\ar[d]\\
\Gamma(A^{\real})\ar[ru]^{\widetilde{\rho}}\ar[r]^{\rho}& \mathfrak{X}(M).}
\end{equation}
Note that, the map $\widetilde{\rho}$ lifts to a CAVB automorphism but not necessarily to a complex almost Lie algebroid automorphism.
We shall use the following identification of vector bundles.
\begin{Lemma}\label{identification_cx}
Let \smash{$N\xhookrightarrow{\iota}M$} be a submanifold. Then, the vector bundles ${\nu(T_{\C}M,T_{\C}N)\!\xrightarrow{\pr_{T_{\C}N}}\! T_{\C}N}$ and \smash{$T_{\C}\nu(M,N)\xrightarrow{T_{\C}\pr_N} T_{\C}N$} are isomorphic as real vector bundles.
\end{Lemma}
\begin{proof}
  We recall that $T_{\C}M=TM\times_M TM$ and so $ T(T_{\C}M)=TTM\times_{TM} TTM$. Consequently,
  \[ T(T_{\C}M)|_{T_{\C}N}=(TTM\times_{TM} TTM)|_{TN\times_{N}TN}=TTM|_{TN}\times_{TM|_N} TTM|_{TN}.\]
  Note that
  \begin{align*}
  \nu(T_{\C}M,T_{\C}N)&{}=(T(T_{\C}M)|_{T_{\C}N})/T(T_{\C}N)\\
  &{}=\bigl(TTM|_{TN}\times_{TM|_N} TTM|_{TN}\bigr)/TTN\times_{TN} TTN\\
  &{}\cong \nu(TM,TN)\times_{\nu(M,N)}\nu(TM,TN)\\
  &{}\cong T_{\C}\nu(M,N).
\tag*{\qed}
\end{align*}
\renewcommand{\qed}{}
\end{proof}

\begin{Remark}%\label{decomposition_normal_cx_anchor}
   Using the identification of Lemma \ref{identification_cx}, we see that if \smash{$N\xhookrightarrow{\iota} M$} is a transversal to $\rho=\rho_1+{\rm i}\rho_2$, then the anchor map is a map of pairs $\rho\colon \bigl(A,\iota^{!}A\bigr)\to (T_{\C}M, T_{\C}N)$, and $\nu(\rho)\colon \nu\bigl(A,\iota^{!}A\bigr)\to \nu(T_{\C}M, T_{\C}N)\cong T_{\C}\nu(M,N)$ is given by
  \[
  \nu(\rho)=\nu(\rho_1)+{\rm i}\nu(\rho_2).
  \]
\end{Remark}
\begin{Lemma}
   If $X$ is an Euler-like vector field along $N$ $($see Definition $\ref{defTN-EL})$, then its complex tangent extension $X_{T_{\C}}$ $($see Definition $\ref{cx_tangent_extension})$ is Euler-like along $T_{\C}N$.
\end{Lemma}

\begin{proof}
   It is easy to see that $X_{T_{\C}}|_{T_{\C}N}=0$. Let $\mathcal{E}$ denote the Euler vector field associated to~${\nu(M,N)\rightarrow N}$. Since the complexified tangent lift of the scalar multiplication is the scalar multiplication in $T_{\C}\nu(M,N)\to T_{\C}N$, we note that $\mathcal{E}_{T_{\C}}$ is an Euler vector field. Using the identification of Lemma \ref{identification_cx}, we see that $\nu(X_{T_{\C}})=\nu(X)_{T_{\C}}=\mathcal{E}_{T_{\C}}$.
\end{proof}

We now apply the techniques of \cite{bursztyn2019splitting} to the case of CAVBs.
\begin{Proposition}\label{splitting_cavb}
Let $(A,\rho)$ be an involutive CAVB with constant real rank and let \smash{$N\xhookrightarrow{\iota} M$} be a submanifold transverse to $\rho^{\real}$. Then
\begin{enumerate}\itemsep=0pt
\item[$1.$] There exists $\widetilde{X}\in \mathfrak{aut}(A,\rho)$ such that $\widetilde{X}|_{\iota^! A}=0$ and $X={(pr_{TM})}_* \widetilde{X}$ is Euler-like along $N$.
\item[$2.$] Each choice of vector field $\widetilde{X}$ as in $(1)$ determines an isomorphism of CAVBs
\[
\widetilde{\psi}\colon\ p^! \iota^! A \to A|_U,
\]
whose base map is the tubular neighborhood embedding $\psi\colon \nu_N\to U\subseteq M$ $($see Definition~$\ref{defTN-EL})$.
\end{enumerate}
\end{Proposition}

\begin{proof}
(1) The proof is similar to \cite[Theorem 3.13]{bursztyn2019splitting}, we just have to work on $T_{\C}M$ instead of~$TM$.

   (2) By Lemma \ref{existence_esp_section}, consider $\epsilon\in \Gamma(A^{\real})$ such that $\rho(\epsilon)=X$ is Euler-like along $N$ and also consider the lifting of equation \eqref{lifting_aut}, $\widetilde{X}\in \aut(A, \rho)$, given by $\widetilde{X}=\widetilde{\rho}(\epsilon)$.
  Note that $\widetilde{X}$ is $\rho$-related to $X_{T_{\C}}$, by the description of $\aut(A,\rho)$.
   Then \smash{$\nu\bigl(\widetilde{X}\bigr)$} is $\nu(\rho)$-related to $\nu(X_{T_{\C}})=\mathcal{E}_{T_{\C}}$, the last one is an Euler vector field and by Lemma \ref{fiberwise_isom}, $\nu(\rho)$ is a fibre-wise isomorphism. Consequently,~\smash{$\widetilde{X}$}~is an~Euler-like vector field along $\iota^! A$.
   Consider $\Phi_t$ and $\widetilde{\Phi}_t$ the flow of $X$ and $\widetilde{X}$, respectively. Since both are Euler-like, the maps $\lambda_t=\Phi_{-\log(t)}$ and \smash{$\widetilde{\lambda_t}=\widetilde{ \Phi}_{-\log(t)}$} are defined at $t=0$. Note that $ \lambda_t\circ \psi=\psi\circ \kappa_t$, in the case of $t=0$ this means that $\lambda_0\circ\psi=p\circ \iota$, where $\psi\colon \nu_N\to U\subseteq M$ is the tubular neighborhood embedding associated to $X$. Recall that the maps $\Tilde{\Phi}_t$ are CLA automorphisms with base $\Phi_t$, so the maps $\Tilde{\lambda}_t$ restricted to $A|_U$ are automorphism of CLA with base $\lambda_t$, for all $t\geq 0$.
   Since $\widetilde{\lambda}_t\circ\widetilde{\lambda}_s=\widetilde{\lambda}_{t+s}$, we have that $\widetilde{\lambda}_t\circ\widetilde{\lambda}_t=\widetilde{\lambda}_{2t}$ and taking limit to $t=0$, we get that $\widetilde{\lambda}^2_0=\widetilde{\lambda}_0$ and so $\widetilde{\lambda}_0$ is a projection. Given that $\Tilde{X}|_{\iota^! A}=0$, we have that $\Tilde{\lambda}_t|_{\iota^! A}={\rm Id}_{\iota^! A}$, so taking limit $\Tilde{\lambda}_0|_{\iota^! A}={\rm Id}_{\iota^! A}$, so $\iota^! A\subseteq \im\bigl(\Tilde{\lambda}_0\bigr)$ and using the fact that
  \begin{equation}\label{eq_inclusion}
    \rho\circ \Tilde{\lambda}_0=T_{\C}\lambda\circ \rho=T_{\C}p\circ T_{\C}\psi^{-1}\circ\rho,
  \end{equation}
  we obtain that $\im\bigl(\widetilde{\lambda}_0\bigr)=\iota^! A$. Consider the map
  \begin{align*} \Psi\colon \ A|_U \to \iota^! A\times T_{\C}\nu_N, \qquad
   a\in A \mapsto \bigl(\widetilde{\lambda}_0(a), T_{\C}\psi^{-1}(\rho(a))\bigr).
   \end{align*}
By equation \eqref{eq_inclusion}, the image of $\Psi$ lies in $p^!\iota^! A$. Finally, note that $\Psi$ is an isomorphism from $A|_U$ to $p^!\iota^! A$, and we define $\widetilde{\psi}=\Psi^{-1}$ as the desired map.
\end{proof}


\begin{Corollary}%\label{loc_splitting_cavb}
Let $(A,\rho)$ be an involutive CAVB and $m\in M$. Assume that $(A,\rho)$ has constant real rank in a neighborhood of $m$. Let \smash{$N\xhookrightarrow{\iota}M$} be a submanifold such that $\Delta|_m\oplus T_m N=T_m M$, and set $P=\Delta{|_m}$. Then, there exists a neighborhood $U$ of $m$ such that $A{|_{U}}$ is isomorphic to~${\iota^! A\times T_{\C}P}$.
\end{Corollary}
\begin{proof}
Since $N$ is completely transversal to $P=\Delta_{|_m}$, we can choose a trivialization of the normal bundle $\nu_N=N\times P$. Then, by Proposition \ref{splitting_cavb}\,(2), we have
$p^!\iota^! A=\iota^! A\times T_{\C}P$.
\end{proof}

\subsection{Local description of CLA}%\label{locstrcla}
Consider a CLA $(A, [\cdot,\cdot],\rho)$ with constant real rank. Since $A$ is a CLA, there exists a natural lift
$\widetilde{\rho}\colon \Gamma(A^{\real})\to \mathfrak{aut}_{CLA}(A)$ of the anchor map $\rho$, as in equation~\eqref{lifting_aut}, now defined by the Lie algebroid bracket.
\begin{Theorem}\label{splitting_CLA}
Let $(A,[\cdot,\cdot],\rho)$ be a CLA with constant real rank, and let \smash{$N\xhookrightarrow{\iota}M$} be a submanifold transverse to $A^{\real}$ $($and so transverse to $A)$. Choose a section $\epsilon\in \Gamma(A^{\real})$, such that $\epsilon|_N=0$ and $\rho(\epsilon)\in \mathfrak{X}(M)$ is Euler-like along $N$. Then, there exists a tubular neighborhood embedding $\psi\colon \nu_N\to U\subseteq M$, depending on the choice of $\epsilon$, and a Lie algebroid isomorphism
\[
\widetilde{\psi}\colon\ p^{!}\iota^{!}A\to A|_U
\]
  covering the tubular neighborhood embedding $\psi\colon \nu_N\to U\subseteq M$
\end{Theorem}
\begin{proof}
   The proof follows the same strategy as that of Proposition \ref{splitting_cavb}, now using the lift described above. Thus, we obtain a map $\widetilde{\psi}\colon p^{!}\iota^{!}A\xrightarrow{}A|_U$ that is an isomorphism as CAVBs, we just have to check that $\widetilde{\psi}$ preserves the Lie algebroid bracket. To this end, we use the argument of \cite[Remark 3.20]{bursztyn2019splitting}. Consider the family of maps
   \begin{align*}
   \Psi_t\colon \ A|_U \to \iota^! A\times T_{\C}\nu_N,\qquad
   a\in A \mapsto \bigl(\widetilde{\lambda}_t(a), T_{\C}\psi^{-1}(\rho(a))\bigr).
   \end{align*}
    Note that the image of $\Psi_t$ is $\kappa_t^! \psi^! A$ and consider the maps \smash{$\widehat{\psi}_t=\Psi^{-1}_t$}. Since each map \smash{$\widehat{\psi}_t$} preserves the bracket for all $t>0$, it follows that their limit \smash{$\widehat{\psi}_0=\widetilde{\psi}$} also preserves it. This completes the~proof.
\end{proof}
